% Seção II — Referencial Teórico
% Adicionar uma \subsection por conceito principal. Mínimo de 3 subseções.
% Cada subseção deve terminar com pelo menos uma referência \cite{}
\section{Referencial Teórico}
\label{sec:referencial}

\subsection{[Conceito Fundamental 1]}
\label{sec:conceito1}

[Definição e contextualização do primeiro conceito fundamental.
Apresentar o estado da arte, definições consagradas na literatura e
sua relevância para o trabalho proposto.
Aproximadamente 2 a 3 parágrafos. \cite{ref1}]

\subsection{[Conceito Fundamental 2]}
\label{sec:conceito2}

[Definição e contextualização do segundo conceito fundamental.
Relacionar com o conceito anterior e mostrar como ambos se complementam
no contexto deste trabalho.
Aproximadamente 2 a 3 parágrafos. \cite{ref2}]

\subsection{Tecnologias e Ferramentas Utilizadas}
\label{sec:tecnologias}

[Descrever as principais tecnologias empregadas no projeto,
fundamentando cada escolha com referências bibliográficas.
Abordar versões, licenças e casos de uso relevantes. \cite{ref3}]
